\chapter{Appendix}
\label{ch::Appendix}

\setcounter{chapter}{1}

\section{Julia programming language initial setup}
\label{sec::JuliaSetup}

This study will be implemented using the Julia programming language.
This section provides an overview of the initial setup required to start programming in Julia.

The Julia programming language is a high-level, high-performance language designed for
numerical and scientific computing. It is known for its speed and ease of use, making
it an ideal choice for developing computational applications. This section provides an
overview of the initial setup required to start programming in Julia.

\subsection{Installation}

The first step in setting up Julia is to download and install the latest version of the language
from the official website (\url{https://julialang.org/downloads/}). The installation process is
straightforward and involves running the installer and following the on-screen instructions.

This will enable you to run Julia code from the command line but to take full advantage, we will use
Visual Studio Code and the Julia extension for Visual Studio Code.

\subsection{Package Management in Julia}

Managing packages efficiently is essential when working on different Julia projects. Julia provides a built-in package manager, \texttt{Pkg}, which allows users to install, update, and manage dependencies. However, to ensure reproducibility and avoid conflicts between projects, it is best to use project-specific environments.

Each Julia project organizes its dependencies using two key files:
\begin{itemize}
    \item \textbf{\texttt{Project.toml}}: Specifies the direct dependencies of the project.
    \item \textbf{\texttt{Manifest.toml}}: Stores the full list of installed packages, including exact versions.
\end{itemize}

Working with project environments ensures that each project has its own set of dependencies,
preventing conflicts between different projects. You will just need the \textbf{\texttt{Project.toml}} file to recreate the environment.


\subsubsection*{Automating Package Management}

To automate package management tasks such as activating environments, resolving
dependencies, and installing required packages, we can use the \texttt{setup.jl} script.

\lstinputlisting[language=Julia, caption=setup.jl]{code/1_annexes/setup.jl}


\section{System Information and Hardware Overview}

It's important to know the system and hardware specifications to know the upper limits of the hardware so that we can compare the results obtained in the benchmarks.

The following code snippets were extracted from the \texttt{system\_info.jl} script, which provides information about the system and hardware specifications.
This script is used to gather information about the system and hardware specifications, such as the operating system, CPU, and GPU details \cite{GPUComputeCapability}.

Function \texttt{cpu\_info(true)} returns the CPU max GFLOPS and function \texttt{gpu\_info(true)} return the GPU max GFLOPS. Using the \texttt{false} argument leds to print the hardware information on the screen.

\begin{lstlisting}[language=Julia, caption=CPU function system\_info.jl]
function cpu_info(return_max_gflops=false)
    N_cores = num_physical_cores()
    N_threads = num_virtual_cores()
    cpuid = cpuinfo()
    string_cpuid = string(cpuid)
    cpu_frequency = 3.7  # GHz (Adjust as needed)
    AVX_value = get_avx_value(string_cpuid)
    M_ops = 4  # 2 FMA, 2 operations each (multiply-add) (AVX2)

    if return_max_gflops
        if AVX_value > 0
            Theoretical_time = 1e9 / (cpu_frequency * 1e9 * AVX_value * M_ops * N_cores)
            max_gflops = 1 / Theoretical_time
            return max_gflops
        else
            return 0
        end
    else
        println("CPU Information")
        println(" ")
        println("CPU Cores: ", N_cores)
        println("CPU Threads: ", N_threads)
        println("CPU Frequency: ", cpu_frequency, " GHz           (GHz adjusted manually)")
        println("AVX support: ", occursin("256 bit", string_cpuid))
        println("AVX-512 support: ", occursin("512 bit", string_cpuid))
        println("AVX Value: ", AVX_value)
        println("M_ops: ", M_ops, " (AVX2: 2 FMA, 2 operations each (multiply-add))")
        if AVX_value > 0
            Theoretical_time = 1e9 / (cpu_frequency * 1e9 * AVX_value * M_ops * N_cores)
            max_gflops = 1 / Theoretical_time
            println("CPU Theoretical max GFLOPS: ", round(max_gflops, digits=2), " GFLOPS (using base frequency)")
        else
            println("AVX not supported. Cannot compute theoretical GFLOPS.")
        end
    end
end
\end{lstlisting}

\begin{lstlisting}[language=Julia, caption=GPU function system\_info.jl]
function gpu_info(return_max_gflops=false)
    device = CUDA.device()
    capability = CUDA.capability(device)
    sm_count = CUDA.attribute(device, CUDA.DEVICE_ATTRIBUTE_MULTIPROCESSOR_COUNT)
    cores_per_sm = cuda_cores_per_sm(capability)
    total_cuda_cores = sm_count * cores_per_sm
    clock_rate = CUDA.attribute(device, CUDA.DEVICE_ATTRIBUTE_CLOCK_RATE) / 1e6

    max_gflops = 2 * clock_rate * total_cuda_cores
    
    memory_clock_khz = CUDA.attribute(device, CUDA.DEVICE_ATTRIBUTE_MEMORY_CLOCK_RATE)
    bus_width_bits = CUDA.attribute(device, CUDA.DEVICE_ATTRIBUTE_GLOBAL_MEMORY_BUS_WIDTH)
    memory_clock_ghz = (memory_clock_khz / 1e6) * 2
    memory_bandwidth = memory_clock_ghz * (bus_width_bits / 8)  # GB/s (8 bits = 1 byte)

    if return_max_gflops
        return max_gflops
    else
        println("GPU Information")
        println(" ")
        println("GPU Name: ", CUDA.name(device))
        println("GPU Compute Capability: ", capability.major, ".", capability.minor, " (", get_GPU_architecture(capability), ")")
        println("GPU Memory: ", round(CUDA.totalmem(device) / 1e9, digits=2), " GB")
        println("GPU Memory: ", round(CUDA.totalmem(device) / 2^30, digits=2), " GiB")
        println("GPU Memory Bandwidth: ", round(memory_bandwidth, digits=2), " GB/s")
        println("GPU Streaming Multiprocessor (SM) Count: ", sm_count)
        println("CUDA Cores per SM: ", cores_per_sm)
        println("Total CUDA Cores: ", total_cuda_cores)
        println("GPU Clock Rate: ", clock_rate, " GHz")
        println("GPU Theoretical max GFLOPS: ", round(max_gflops, digits=2), " GFLOPS")
    end
end
\end{lstlisting}

It's important to modify the \texttt{system\_info.jl} script to include any additional hardware information you want to display and customize parameters as the \texttt{cpu\_frequency} or \texttt{M\_ops}.

\clearpage